\begin{frame}{PSUM = TLE}
    \metroset{block=fill}

    Vamos cortar a dependência de $p_3$ para $p_2$ (isto é, cortar a aresta $p_3 \to p_2$). Neste caso, teremos $p_1 = v_1$ e $p_3 = v_3$, então desconsideraremos $p_1$ e $p_3$. Temos:
    \begin{center}
        \tikz {
        \node (1) at (0,0) {$v_1$};
        \node (2) at (1,0) {$v_2$};
        \node (3) at (2,0) {$v_3$};
        \node (4) at (3,0) {$v_4$};

        % \node (p1) at (0,1) {$p_1$};
        \node (p2) at (1,1.5) {$p_2$};
        % \node (p3) at (2,2) {$p_3$};
        \node (p4) at (3,2.5) {$p_4$};

        \graph { 
            (p4) -> {(4), (3)};
            % (p3) -> {(3), (p2)};
            (p2) -> {(2), (1)};
            % (p1) -> (1);
        };
        }
    \end{center}
    
    Temos um problema: como agora $p_4 = v_3+v_4$, não temos uma variável que tenha o valor $v_1+v_2+v_3+v_4$. Adicionamos um valor $r = p_2 + p_4$ e reorganizamos a árvore:

    % \begin{block}{Implementação da busca binária}
    %     \footnotesize
    %     \inputminted{cpp}{codigos/busca.binaria.cpp}    
    % \end{block}

\end{frame}
